Throughout Sections~\ref{sec:shortrun} and~\ref{sec:results} we held the set of steps assigned to the worker fixed, so that AI could improve only the time those steps took to execute.
Here we relax this assumption and allow the firm to restructure production fully around the capabilities of AI.

Over longer horizons the firm can readjust its organizational structure, redrawing the boundaries between jobs and with them the allocation of steps to workers, while optimizing its use of AI.
Moreover, adjustments in the labor market let the firm pay workers according to the actual skill requirements of the jobs it designs, rather than at the rate the short run held fixed, effectively endogenizing wages.
Below we develop the firm's long-run production problem.

\subsection{Jobs and Wages}
\label{sec:longrun.jobdesign}

Just as an AI strategy partitions the steps into tasks, the firm's long-run choice partitions those tasks into jobs.
Definition~\ref{def:job_design} formalizes that second layer of bundling, which Figure~\ref{fig:hierarchical_production} adds to the example of Figure~\ref{fig:steps_tasks}.

\begin{definition}[Job Design]
\label{def:job_design}
A \emph{job} is a contiguous block of tasks $J = (T_b, \dotsc, T_{b'})$, with $b \leq b'$, assigned to a single worker.
A \emph{job design} is a partition $\mathcal{J} = (J_1, \dotsc, J_p)$ of the task sequence $\T$ into $p$ such jobs, one for each worker the firm hires.
\end{definition}
In the bottom two layers of Figure~\ref{fig:hierarchical_production}, Tasks~1 and~2 form Job~1, Task~3 forms Job~2, and Tasks~4 and~5 form Job~3.
Notice that a worker's job can potentially contain both manual and AI chain tasks.

\begin{figure}[!t]
  \centering
  \caption{Illustrative Example of a Firm's Production in the Long Run}
  \label{fig:hierarchical_production}
  % Hierarchical Production Visualization
\begin{adjustbox}{center,trim=0pt 0pt 0pt 0pt,clip,margin=-2cm 0cm}
  \begin{tikzpicture}[scale=0.73, transform shape,
    node distance=0.35cm and 0.91cm,
    every node/.style={rectangle, rounded corners, draw, align=center, minimum width=2cm, minimum height=1cm},
    manual/.style={fill=gray!10},
    automated/.style={fill=green!30},
    augmented/.style={fill=orange!30},
    dashedbox/.style={draw, rectangle, rounded corners, dashed, inner sep=0.25cm, line width=1.25pt},
    humanbox/.style={rectangle, rounded corners, draw, align=center, minimum width=2.25cm, minimum height=1cm, fill=olive!20}, % changed to olive
    aibox/.style={rectangle, rounded corners, draw, align=center, minimum width=2.25cm, minimum height=1cm, fill=blue!20},
    jobbox/.style={rectangle, rounded corners, draw, align=center, minimum width=2.25cm, minimum height=1cm, fill=cyan!20},
    industrybox/.style={rectangle, rounded corners, draw, align=center, minimum width=3cm, minimum height=1cm, fill=teal!60},
    >=latex
  ]

    % Top layer (Steps)
    \node[manual] (S1) {Step 1\\(Manual)};
    \node[automated, right=of S1] (S2) {Step 2\\(Automated)};
    \node[automated, right=of S2] (S3) {Step 3\\(Automated)};
    \node[augmented, right=of S3] (S4) {Step 4\\(Augmented)};
    \node[manual, right=of S4] (S5) {Step 5\\(Manual)};
    \node[augmented, right=of S5] (S6) {Step 6\\(Augmented)};
    \node[manual, right=of S6] (S7) {Step 7\\(Manual)};

    % Arrows between steps (specific arrows as curly lines)
    \draw[->, thick] (S1) -- (S2);
    \draw[->, thick] (S2) -- (S3);
    \draw[->, thick] (S3) -- (S4);
    \draw[->, thick, red, decorate, decoration={snake, amplitude=.5mm, segment length=3.05mm}] (S4) -- (S5);
    \draw[->, thick, red, decorate, decoration={snake, amplitude=.5mm, segment length=3.05mm}] (S5) -- (S6);
    \draw[->, thick] (S6) -- (S7);

    % Dashed boxes (top layer, colors updated)
    \node[dashedbox, draw=olive, fit=(S1)] (DS1) {};
    \node[dashedbox, draw=blue, fit=(S2)(S3)(S4)] (DS2-4) {};
    \node[dashedbox, draw=olive, fit=(S5)] (DS5) {};
    \node[dashedbox, draw=blue, fit=(S6)] (DS6) {};
    \node[dashedbox, draw=olive, fit=(S7)] (DS7) {};

    \path (S1.west) -- (S7.east) coordinate[midway] (CenterTop);

    % Middle layer (Tasks, human task color updated)
    \node[humanbox, below=3.2cm of CenterTop, xshift=-6cm] (T1) {Task 1\\(Human)};
    \node[aibox, right=of T1] (T2) {Task 2\\(AI Chain)};
    \node[humanbox, right=of T2] (T3) {Task 3\\(Human)};
    \node[aibox, right=of T3] (T4) {Task 4\\(AI Chain)};
    \node[humanbox, right=of T4] (T5) {Task 5\\(Human)};

    % Arrows between tasks (specific arrows as curly lines)
    \draw[->, thick] (T1) -- (T2);
    \draw[->, thick, red, decorate, decoration={snake, amplitude=.5mm, segment length=3.05mm}] (T2) -- (T3);
    \draw[->, thick, red, decorate, decoration={snake, amplitude=.5mm, segment length=3.05mm}] (T3) -- (T4);
    \draw[->, thick] (T4) -- (T5);

    % Dashed boxes (middle layer)
    \node[dashedbox, draw=cyan, fit=(T1)(T2)] (DT1-2) {};
    \node[dashedbox, draw=cyan, fit=(T3)] (DT3) {};
    \node[dashedbox, draw=cyan, fit=(T4)(T5)] (DT4-5) {};

    % Job layer nodes
    \node[jobbox, below=2.25cm of T3] (J2) {Job 2};
    \node[jobbox, left=of J2] (J1) {Job 1};
    \node[jobbox, right=of J2] (J3) {Job 3};

    % Arrows between jobs (specific arrows as curly lines)
    \draw[->, thick, red, decorate, decoration={snake, amplitude=.5mm, segment length=3.05mm}] (J1) -- (J2);
    \draw[->, thick, red, decorate, decoration={snake, amplitude=.5mm, segment length=3.05mm}] (J2) -- (J3);

    % Arrows connecting layers (unchanged)
    \draw[->, thick, cyan] (DT1-2.south) -- (J1.north);
    \draw[->, thick, cyan] (DT3.south) -- (J2.north);
    \draw[->, thick, cyan] (DT4-5.south) -- (J3.north);

    \draw[->, thick, olive] (DS1.south) -- (T1.north);
    \draw[->, thick, blue] (DS2-4.south) -- (T2.north);
    \draw[->, thick, olive] (DS5.south) -- (T3.north);
    \draw[->, thick, blue] (DS6.south) -- (T4.north);
    \draw[->, thick, olive] (DS7.south) -- (T5.north);

  \end{tikzpicture}
\end{adjustbox}

  \par\vspace{\fignotesskip}
\begin{minipage}{1\linewidth}
\begin{spacing}{0.2}
{\fontsize{9.5pt}{10pt}\selectfont Notes: The top two layers reproduce Figure~\ref{fig:steps_tasks}, with $m=7$ steps grouped by dashed boxes into the $n=5$ tasks below them.
This figure adds the third layer: cyan dashed boxes group tasks into $p=3$ jobs, each assigned to a separate worker.
Red curly arrows mark the hand-offs incurred at job boundaries, and can be traced from the bottom layer up to the tasks and steps they separate.}
\end{spacing}
\end{minipage}
\end{figure}

Completing a step demands not only time but also skill, by which we mean the body of knowledge and the specific know-how a worker must acquire in order to perform it.
The short run had no need of this notion, as we assumed workers had acquired the skills necessary to perform all activities in their job.
In the long run, however, the firm chooses which steps to assign to which worker, effectively changing the skill requirement of their jobs.
We thus introduce skills to capture this other work dimension.

Accordingly, we extend Definitions~\ref{def:manual_step}--\ref{def:ai_chain} to incorporate the skill requirement of steps as well: 
the skill and time requirement of performing step $i$ manually is $(\manualSkill{i},\, \manualTime{i})$, while verifying the output of one AI attempt at it costs $(\AIskill{i},\, \AItime{i})$.
The two values need not coincide, since doing the work oneself calls for knowing how to do it, whereas verifying what an AI returns may call for knowledge of a different kind.
As with time, a task inherits the skill of whichever mode executes it, so each task $T_b$ in the task sequence $\T = (T_1, \dotsc, T_n)$ carries a pair of skill and time requirements $(\skillcost{b}, \timecost{b})$: it is $(\manualSkill{i},\, \manualTime{i})$ if $T_b$ is the manual step $s_i$, and $\bigl(\AIskill{r},\, \AItime{r}/\prod_{i=\ell}^{r} q_i\bigr)$ if $T_b$ is an AI chain spanning $(s_\ell, \dotsc, s_r)$.
Both cost components of a chain's pair (for a single verification attempt) are therefore inherited from its augmented endpoint.
In particular, the chain demands only the skill required to verify the augmented step's output, and an automated step inside the chain demands no skill of anyone, since no human ever handles it directly.

The firm hires one worker for each job in its job design $\mathcal{J}$.
The worker for job $J_j$ completes all tasks associated with their job.
The total time needed to complete all tasks in job $J_j$ is $\sum_{T_b \in J_j} \timecost{b}$.
The total skill required to complete the tasks in job $J_j$ is $\sum_{T_b \in J_j} \skillcost{b}$.\footnote{
We assume that the skills different steps require are non-overlapping, and thus add up linearly in jobs.
Human capital requirements of steps do overlap in practice, but modeling that overlap would burden the results this section derives with complexities that obscure the main channels we aim to describe.
That said, the results rest on a considerably weaker version of this assumption, namely that each step demand some increment of skill no other step supplies, so that a worker who takes on more steps in their job must acquire more skill.
Any strictly increasing functional form delivers this; linear addition, which we adopt here, is simply the most tractable.

}
We then assume that a worker's wage is proportional to the skill level of their job.\footnote{
One motivation for this formulation is workers paying a human capital investment to acquire the skills necessary for a job.
This investment must be offset by wages, and we express skill levels in units of their corresponding requisite wage.
Rather than formalizing such a wage model here, we impose this as an assumption and provide a more detailed wage formulation in Appendix~\ref{sec:aggregation}.
}
That is,
\begin{align}
\label{eq:wage}
\text{Wage}_j = \sum_{T_b \in J_j} \skillcost{b}.
\end{align}

The firm must pay workers their wage per unit of time regardless of which tasks they are assigned to complete at any given moment.
%\footnote{
%This derivation implicitly assumes a common base wage rate; i.e., all wage differentiation is driven by the human capital cost of acquiring skills.
%More generally, different tasks $T_b$ within job $J_j$ could have different base wage rates $w_b$ that could be influenced, for example, by the supply of and demand for labor of the corresponding type.
%One special case, explicitly motivated and discussed in Appendix~\ref{sec:aggregation}, assumes manual tasks are executed by manual labor at base wage rate $w_{\manualLetter}$, and AI chains are executed by AI management labor at base wage rate $w_{\AIletter}$.
%The formulation in \eqref{eq:wage} implicitly normalizes these base wage rates to $1$, incorporating them into skill costs $\skillcost{b}$ for notational simplicity.
%We maintain this normalization throughout subsequent sections until we explicitly distinguish human labor and AI management labor base wage rates in Appendix~\ref{sec:aggregation}.
%\label{foot:wage}
%}
Taking into account both time and skill requirements, the total wage bill paid to a worker assigned to job $J_j$, per unit of output, is
\begin{align}
\label{eq:jobcost}
\text{WageBill}_j = \Biggl(\sum_{T_b \in J_j} \skillcost{b} \Biggr)\Biggl(\sum_{T_b \in J_j} \timecost{b} \Biggr).
\end{align}
Note what Equation~\eqref{eq:jobcost} implies for a job.
The worker is compensated for every skill the job demands over its whole duration, even though at any given moment they are performing a single task and drawing on only the portion of their knowledge that task requires.
The remainder of their human capital sits idle at that moment and is paid for regardless.

A university professor is a concrete example.
The job of a professor bundles research and teaching together with the small chores that come with them, and the wage it commands reflects the skill the bundle as a whole demands.
Wiping the board demands almost none of that skill, yet the minutes spent on it are paid at the professor's rate rather than at what the chore would command on its own.
For those minutes the expertise that makes the professor worth hiring sits unused, and the university pays for it all the same.

\subsection{Coordination and Hand-Off Cost}
\label{sec:longrun.handoff}

Absent any frictions, the firm would minimize production costs by specializing as far as it can.
Equation~\eqref{eq:jobcost} implies dividing a job in two splits its skill requirement between two workers while leaving the total time to be worked unchanged, so it strictly lowers the wage bill.
Carried to its conclusion, this logic would have the firm fully specialize work and assign every task to its own worker, each of whom acquires exactly the knowledge their task requires and none of whom is ever paid for skill they are not using.

In reality, work is not organized this way, because dividing production among many workers is costly.
Each division creates a point at which work must pass from one person to the next, and coordinating that intermediate output takes time: whoever receives it has to be told what was done, why, and what state the output is now in, and whoever hands it over has to put it in a form that can be received.\footnote{\cite{lee2023task} formalize such interdependencies and construct an occupation-level index of coordination intensity from the O*NET dataset, which they relate to wages and to exposure to AI.}
We model this coordination friction with an additional \emph{hand-off} cost, in time, that a worker spends passing their work on to the next worker, incurred whenever $J_j$ is not the final job in the job design $\mathcal{J}$.

Because production in our model exhibits an order, we can be more specific about what a hand-off costs.
A hand-off happens at a particular point in production, and what it costs depends on which point.
Some places are natural seams, where the work has reached a self-contained intermediate output, a finished design or a signed-off specification, that can be described briefly and passed on.
Others fall in the middle of something, where whoever receives the work has to reconstruct the context the previous worker had built up.
We therefore attach the hand-off cost to the step at which the break occurs rather than to the worker or to the number of jobs, so that \emph{where} the firm divides the workflow matters and not merely how often it divides.
The red curly arrows in Figure~\ref{fig:hierarchical_production} illustrate that hand-offs occur at job boundaries and that their cost reflects the characteristics of the last step in a worker's job.

Formally, given step $s_i$, we write $\handofftime{i}$ for the additional hand-off time spent by a worker for whom the last step of their job is $s_i$, and $\handofftime{}(J_j)$ for the hand-off time of job $J_j$, which is equal to $\handofftime{i}$ if $s_i$ is the final step in $J_j$.
The final job hands off to no one, so $\handofftime{m} = 0$.
Conditional on where the hand-off occurs in the production process, its cost does not otherwise depend on previous hand-off events.
Thus the total time needed to complete job $J_j$, including hand-off costs, is
\[
\handofftime{}(J_j) + \sum_{T_b \in J_j} \timecost{b}.
\]

Because the worker is paid for as long as their job takes, this time enters the wage bill exactly as execution time does, so Equation~\eqref{eq:jobcost} becomes
\begin{align}
\label{eq:wagebill}
\text{WageBill}_j = \Biggl(\sum_{T_b \in J_j} \skillcost{b} \Biggr) \Biggl(\handofftime{}(J_j) + \sum_{T_b \in J_j} \timecost{b} \Biggr).
\end{align}
Equation~\eqref{eq:wagebill} highlights the two opposing forces that shape a job's boundaries.
Adding more tasks to a worker's job increases the cumulative skill requirement and thus the wage rate for that worker.
However, if tasks are kept separate, workers must incur additional hand-off costs at each job boundary.
These two effects jointly determine the optimal degree of task aggregation.

Going back to the university professor example, this is perhaps why the chores fall to the professor in the first place. On wages alone the university should hand the cleaning of the board to a cheaper worker. What stops it is that the work would have to change hands mid-lecture: the professor must find someone, the lecture waits, and the worker must be told what is wanted. That costs more than the professorial minutes it saves, so the professor's job absorbs the chore.

We explore this trade-off in greater detail in Section~\ref{sec:longrun.specialization}, after introducing the problem the firm solves in choosing its jobs.


\subsection{Long-Run Production}
\label{sec:longrun.production}

Given the sequence of steps $\mathcal{S} = (s_1, \dotsc, s_m)$, the firm can design both the partition $\T$ of steps into tasks and the partition $\mathcal{J}$ of tasks into jobs.
The choice of $\T$ determines the time and skill requirements of each task.
Given $\T$, the choice of $\mathcal{J}$ then determines worker wage rates and time needed per unit of output (including hand-off costs).
Formally, we can write $\mathcal{P}(X)$ for the set of partitions of a sequence $X$ into contiguous subsequences.
Then the full optimization problem faced by the firm can be expressed as
\begin{align}
\label{eq:totalcost_with_handoff}
\min_{\T \in \mathcal{P}(\mathcal{S})} \
\min_{\mathcal{J} \in \mathcal{P}(\T)} \
\text{TotalCost}(\mathcal{J}; \T)
= \sum_{J_j \in \mathcal{J}} \text{WageBill}_j =
\sum_{J_j \in \mathcal{J}} \Biggl[\Biggl(\sum_{T_b \in J_j} \skillcost{b} \Biggr) \Biggl(\handofftime{}(J_j) + \sum_{T_b \in J_j} \timecost{b} \Biggr)\Biggr]
\end{align}
where recall that, for each $T_b \in \T$, $\timecost{b}$ is as described in Section~\ref{sec:model.tasks} and $\skillcost{b}$ as above, both depending on the selected mode of operation for individual steps.

We refer to \eqref{eq:totalcost_with_handoff} as the firm's long-run optimization problem, as it anticipates the adjustment of worker wages to fit the skill requirements of each job and sets job responsibilities accordingly.
It nests the short-run problem in \eqref{eq:totalcost_shortterm} of Section~\ref{sec:shortrun.production} as the case of a single job spanning the whole production, with wage being fixed rather than chosen.
%In Section~\ref{sec:extensions.computation} we discuss how the solution to Problem~\eqref{eq:totalcost_with_handoff} can be approximated to any desired accuracy in polynomial time.

Next we discuss the implications of this optimization problem.


\subsection{AI and the Limits of Worker Specialization}
\label{sec:longrun.specialization}

Here we discuss how AI affects the boundaries of jobs.
Because the firm chooses those boundaries when it solves Problem~\eqref{eq:totalcost_with_handoff}, the answer runs through the trade-off that fixes them, between the skill a broad job demands and the coordination a narrow one requires.
We set that trade-off out first, and then ask how AI moves it.

The wage bill in \eqref{eq:wagebill} is a product of a skill and a time, so it reads naturally as an area, which is what makes the trade-off visible.
Draw each task as a rectangle whose height is its skill cost $\skillcost{b}$ and whose width its time cost $\timecost{b}$; a job's wage bill is then the area of the rectangle that spans the tasks it contains, as tall as the total skill it demands and as wide as the total time it takes.
Figure~\ref{fig:lr_job_design} draws the two-task case, the smallest in which the firm has a choice to make.
\begin{figure}[!t]
    \centering
    \caption{The Trade-off between Worker Specialization and Coordination}
    \label{fig:lr_job_design}
    \includegraphics[width=1\linewidth]{plots/job_design.png}
    \par\vspace{\fignotesskip}
\begin{minipage}{1\linewidth}
\begin{spacing}{0.2}
{\fontsize{9.5pt}{10pt}\selectfont Notes: The blue bounded rectangles are tasks, with height the task's skill requirement ($\skillcost{}$) and width its time requirement ($\timecost{}$).
The shaded areas are the wage bills of jobs.
In the left panel the two tasks are assigned to two specialized workers, and the pink rectangle is the hand-off cost incurred between them.
In the right panel they are bundled into a single job, which eliminates the hand-off but requires a worker skilled enough to perform both tasks and paid at the corresponding rate throughout.}
\end{spacing}
\end{minipage}
\end{figure}
In the left panel the tasks go to two specialized workers, each of whom acquires only the skill their own task needs, so no one is paid for knowledge they are not using; but the work must pass between them, and the pink rectangle is what that hand-off costs.
Notably, the hand-off shares the skill requirements of the worker doing it for an extra amount of time added to the last task in their job.
In the right panel one worker performs both tasks, which removes the hand-off but obliges that worker to hold both skills throughout, the dead area sitting beside whichever task they are not currently performing.
Bundling therefore buys the removal of a hand-off at the price of idle skill, and where those two prices balance determine the firm's degree of specialization.

With more than two tasks the optimum need not sit at either extreme.
Consider a workflow that, under a fixed AI deployment strategy $\T$, has $n=3$ tasks with the cost parameters
\[ \bigl(\skillcost{b},\, \timecost{b},\, \handofftime{b}\bigr)_{b=1,2,3} = (3,\,1,\,3),\quad (1,\,2,\,0.5),\quad (2,\,2,\,0). \]
With three tasks there are four ways the firm can design jobs,\footnote{
More generally, a production sequence with $n$ tasks has $2^{(n-1)}$ unique job designs.
} and Figure~\ref{fig:job_design} draws all four in the same terms, once with hand-off costs switched off and once with them on; Table~\ref{tab:job_design} in Appendix~\ref{app:tables_and_figures} reports the cost of each job design.
\begin{figure}[!t]
\begin{center}
\caption{Geometric Interpretation of Optimal Job Design} \label{fig:job_design}
\begin{subfigure}[b]{0.49\textwidth}
    \caption{Job Design Without Hand-off Costs}
    \includegraphics[width=\textwidth]{plots/combined_grid_no_handoff.png}
\end{subfigure}
\begin{subfigure}[b]{0.49\textwidth}
    \caption{Job Design With Hand-off Costs}
    \includegraphics[width=\textwidth]{plots/combined_grid_with_handoff.png}
\end{subfigure}
\end{center}
\par\vspace{\fignotesskip}
\begin{minipage}{1\linewidth}
\begin{spacing}{0.2}
{\fontsize{9.5pt}{10pt}\selectfont Notes: Each panel draws all four job designs of the three-task workflow in the convention of Figure~\ref{fig:lr_job_design}, with the area of each job's bounding box its wage bill. The cost-minimizing design is outlined in red. Panel~(a) sets hand-off costs to zero, in which case full specialization is optimal. Panel~(b) introduces them, and the optimum moves to $[1,2][3]$. The skill, time, and hand-off costs of the three tasks are given in the text, and Table~\ref{tab:job_design} in Appendix~\ref{app:tables_and_figures} reports the cost of each design.}
\end{spacing}
\end{minipage}
\end{figure}
Panel~(a) sets hand-off costs to zero and recovers the frictionless benchmark: the cost-minimizing design, outlined in red, is full specialization, with each task its own separate job.
Panel~(b) reinstates hand-offs and the optimum moves to $[1,2][3]$, in which tasks 1 and 2 share a job and task 3 stands alone.
The reason is the large hand-off time $\handofftime{1}=3$ that a boundary between tasks 1 and 2 would trigger: bundling them avoids that hand-off at the price of a worker who must hold both skills and be paid for both throughout.

This is the machinery AI acts on, and it does so through two channels.
The first is skill. An AI chain leaves the worker holding only the skill needed to verify its endpoint, so the tasks AI absorbs become short rectangles in the picture above.
That cuts the idle skill a bundled job pays for, which makes broader jobs cheaper to sustain and pushes the firm toward consolidation.
But the same reduction works in the opposite direction as well, because what a hand-off costs is also a skill multiplied by a time: a worker whose job AI has made less skill-intensive is cheaper to leave idle at a boundary, so the hand-off the firm was bundling to avoid becomes less worth avoiding.
In the workflow above, for instance, tasks 1 and 2 share a job because a skill-$3$ worker would otherwise sit through a hand-off of length $3$; were AI to take over task 1 and leave only verification skill behind, that reason would weaken along with the wage rate it applies to.

The second channel is time.
For many steps, checking what AI produces takes less time (in expectation) than performing them, so the tasks AI absorbs become narrow rectangles as well as short ones.
Narrower tasks cut the idle skill a bundled job pays for, since what bundling wastes is one task's skill sitting through another task's time, and that pushes toward consolidation as the skill reduction did.
The hand-off answers to the same force.
We have treated $\handofftime{i}$ as a primitive to capture the tacit and social knowledge of transferring work, but AI plausibly lowers it as well, by absorbing the peripheral coordination a transfer requires, such as summarizing what was done or drafting the documentation that accompanies it.
Cheaper hand-offs make every boundary cheaper to draw, which pushes toward narrower jobs and more specialization.
Panel~(a) of Figure~\ref{fig:job_design} is the limiting case, with hand-offs free and full specialization optimal.
Time therefore divides as skill does: a shorter task makes bundling cheaper, a shorter hand-off makes splitting cheaper, and the two work against each other.

What this framework delivers is therefore not a prediction about job breadth but an account of what determines it.
AI presses on both sides of the specialization trade-off at once, and which way jobs move depends on the relative size of the skill and time it strips out and the coordination it leaves behind.
That is an empirical question, and one the framework poses in a form that says what would have to be measured to answer it.